% Original source preamble and source-level macro declarations.
% This archival file is not the blueprint driver preamble.

\documentclass[10pt]{amsart}
\usepackage[utf8]{inputenc}


\usepackage{amsmath}
\usepackage{amsfonts}
\usepackage{amssymb}
\usepackage{amsthm}
\usepackage{mathrsfs}
\usepackage{dsfont}
\usepackage{chngcntr}
\usepackage{mathtools}
\usepackage[all,cmtip]{xy}
\usepackage{tikz}
\usepackage[colorlinks]{hyperref}
\hypersetup{linkcolor=blue,citecolor=green}
\usepackage{cite}
\usepackage{fancyhdr}
\usepackage{stmaryrd}
\usepackage{yfonts}
\usepackage{adjustbox}

\tikzset{double line with arrow/.style args={#1,#2}{decorate,decoration={markings,%
mark=at position 0 with {\coordinate (ta-base-1) at (0,1pt);\coordinate (ta-base-2) at (0,-1pt);},
mark=at position 1 with {\draw[#1] (ta-base-1) -- (0,1pt);
\draw[#2] (ta-base-2) -- (0,-1pt);
}}}}

\newcommand{\veq}{\mathrel{\rotatebox{90}{$=$}}}
\newcommand{\vneq}{\mathrel{\rotatebox{90}{$\neq$}}}



\font\wncyr=wncyr9.8
\newcommand{\sha}{\text{\wncyr{W}}}

\usepackage[top=2.6cm, bottom=2.4cm, left=2.4cm, right=2.4cm]{geometry}

% comm diagram
\usepackage{tikz}
\usepackage{tikz-cd}

% hyperlinks

\numberwithin{equation}{section}




% Theorems
\newtheorem{theorem}{Theorem}
\newtheorem{lemma}[theorem]{Lemma}
\newtheorem{prop}[theorem]{Proposition}
\newtheorem{cor}[theorem]{Corollary}
\newtheorem{conj}[theorem]{Conjecture}
\newtheorem{conv}[theorem]{Convention}
\newtheorem{assu}[theorem]{Assumption}
\newtheorem{notation}[theorem]{Notation}
\newtheorem{ques}[theorem]{Question}

\theoremstyle{definition}
\newtheorem{definition}[theorem]{Definition}
\newtheorem{remark}[theorem]{Remark}
\newtheorem{question}[theorem]{Question}
\newtheorem{example}[theorem]{Example}

\theoremstyle{remark}
\newtheorem*{rem}{Remark}


% Theorems for intro

\newtheorem{theoremintro}{Theorem}

% Commands

\newcommand{\Z}{\mathbb{Z}}

%\usepackage{tcolorbox}


\newcommand{\rightarrowdbl}{\rightarrow\mathrel{\mkern-14mu}\rightarrow}

\newcommand{\xrightarrowdbl}[2][]{%
  \xrightarrow[#1]{#2}\mathrel{\mkern-14mu}\rightarrow}
  
\DeclareRobustCommand{\qed}{%
  \ifmmode
    \eqno \def\@badmath{$$}%$$
    \let\eqno\relax \let\leqno\relax \let\veqno\relax
    \hbox{\openbox}%
  \else
    \leavevmode\unskip\penalty9999 \hbox{}\nobreak\hfill
    \quad\hbox{\openbox}%
  \fi
}

\setcounter{tocdepth}{2}


% Theorems
\renewcommand{\thetheorem}{\arabic{section}.\arabic{subsection}.\arabic{theorem}}
\renewcommand{\thetheoremintro}{\Alph{theoremintro}}


\newcommand{\cL}{\mathscr{L}}
\newcommand{\cN}{\mathscr{N}}

\newcommand{\Lcal}{\mathcal{L}}
\newcommand{\Pp}{\mathfrak{P}}
\newcommand{\Fcal}{\mathcal{F}}
\newcommand{\cF}{\mathcal{F}}
\newcommand{\bZ}{\mathbb{Z}}
\newcommand{\bQ}{\mathbb{Q}}
\newcommand{\Q}{\mathbb{Q}}
\newcommand{\C}{\mathbb{C}}
\newcommand{\rH}{\mathrm{H}}
\newcommand{\res}{k}
\newcommand{\fm}{\mathfrak{m}}
\DeclareMathOperator{\Frob}{Frob}
\DeclareMathOperator{\sel}{Sel}
\DeclareMathOperator{\Hom}{Hom}
\DeclareMathOperator{\ch}{char}
\DeclareMathOperator{\len}{length}
\DeclareMathOperator{\coker}{coker}
\DeclareMathOperator{\loc}{loc}
\DeclareMathOperator{\ord}{ord}
\DeclareMathOperator{\un}{un}
\DeclareMathOperator{\tr}{tr}
\DeclareMathOperator{\im}{im}
\DeclareMathOperator{\GL}{GL}
\DeclareMathOperator{\rk}{rk}
\DeclareMathOperator{\Gal}{Gal}
\DeclareMathOperator{\Aut}{Aut}
\DeclareMathOperator{\End}{End}


\newcommand{\CF}{\mathcal{F}}
\newcommand{\CL}{\mathscr{L}}
\newcommand{\CN}{\mathscr{N}}
\newcommand{\kap}{\kappa}
\newcommand{\ind}{\mathrm{ind}}
\newcommand{\rank}{\mathrm{rank}}


%\setcounter{section}{-1}

% reset counter at each section
\makeatletter
\@addtoreset{theorem}{subsection}
\makeatother


